\documentclass[11pt]{article}
\usepackage[margin=1in]{geometry}
\usepackage{amsmath,amssymb,amsfonts,mathtools}
\usepackage{array,booktabs,enumitem,graphicx,xcolor,hyperref}
\usepackage[T1]{fontenc}
\usepackage[utf8]{inputenc}
\title{A \$U(N)\$ Gauge Theory in Three Dimensions as an Ensemble of Surfaces}
\author{Source-backed arXiv sample}
\date{}
\begin{document}
\maketitle
\begin{abstract}
A particular \$U(N)\$ gauge theory defined on the three dimensional dodecahedral lattice is shown to correspond to a model of oriented self-avoiding surfaces. Using large \$N\$ reduction it is argued that the model is partially soluble in the planar limit.
\end{abstract}
\section{Source Metadata}
This sample is based on arXiv record \texttt{hep-th/9108015}. Authors: F. David, H. Neuberger.
\section{Extracted Prose 1}
RU--91--34 July 29, 1991 1.5cm A math expression Gauge Theory in Three Dimensions as an Ensemble of Surfaces 1cm by 1cm Fran cois David Physique Th\textbackslash{}' eorique CNRS. .2cm Service de Physique Th\textbackslash{}' eorique Laboratoire de l'Institut de Recherche Fondamentale du Commissariat \textbackslash{}` a l'Energie Atomique. de Saclay F-91191 Gif-sur-Yvette Cedex .7cm and .7cm
\section{Extracted Prose 2}
1cm A particular math expression gauge theory defined on the three dimensional dodecahedral lattice is shown to correspond to a model of oriented self-avoiding surfaces. Using large math expression reduction it is argued that the model is partially soluble in the planar limit. .3cm
\end{document}
